% =====================================================================
% Experiments section. All numbers are from logs/paper/benchmark_td3/report.md
% (python -m evaluation.benchmark --episodes 10 --out logs/paper/benchmark_td3),
% figures from figures/exp_results.ipynb.
% Requires booktabs and graphicx. Cite key xu2016modeling must exist in the .bib.
% =====================================================================

\section{Experimental Evaluation}
\label{sec:exp}

% ---------------------------------------------------------------------
\subsection{Setup}
\label{sec:exp:setup}

\paragraph{Simulator}
All controllers are evaluated in a Python simulator of the rover whose power, throughput and
battery models are fitted to measurements. Every method drives the same missions under the
same workload, so the methods are compared on identical conditions.
The motor power is a steady-speed fit to robot measurements~\cite{naseri2025co} with a
temperature-dependent copper loss. The rover runs two perception applications,
\texttt{APP-DT} and \texttt{APP-OP}, each with a throughput requirement of
$q_{\mathrm{ref}} = 25$~event batches per second. Their throughput and the power of the
computing board are fitted to measurements of the on-board NVIDIA Jetson TX2 at its 11 \ac{CPU}
\ac{DVFS} levels (0.50--2.04~GHz) and at three environmental complexity levels
$c \in \{1, 2, 3\}$ (low, medium, high), which differ in the average number of events per
batch~\cite{shahsavari2024coordinated}. The throughput model fits the measured curves with
a root-mean-square error of 0.43~batch/s. The speed range is 0.5--5~m/s, and the control
period is 1~s.

\paragraph{Battery model}
The rover carries a 4S1P 5~Ah LiPo pack (14.8~V nominal). The cell models are identified on
the LiPo aging dataset~\cite{galeotti2023lipo}, using cells 1 and 2 for fitting and cell 3 for
validation, and are scaled to the pack (resistances by $N_s/N_p = 0.81$).
The polynomials $p_j$, the aging coefficients $\kappa_j$ of \eqref{eq:ecm_param} and the two time
constants of the \ac{ECM} (Table~\ref{tb:battery_parameter}) are identified in the time domain from
202 current pulses of 360~s at 1~C, recorded at different aging stages. On the 91 pulses of
the validation cell, the median voltage error is 12~mV.
For the degradation model \eqref{eq:battery_soh_define}--\eqref{eq:xu_cal},
because all cells of the dataset were cycled under a single protocol, the
stress-factor shape ($k_\sigma$, $k_T$, $k_{\delta_1}$--$k_{\delta_3}$, $k_t$) cannot be
identified from it and is taken from~\cite{xu2016modeling}; only the cycle-rate scale
$k_{\mathrm{cyc}}$ is fitted, so that one cycle at the dataset protocol (median depth of
discharge 87\%) reproduces the measured capacity fade of $4.8 \times 10^{-4}$ per cycle.

\begin{table}[t]
\centering
\caption{Battery model parameters (4S1P pack, 5~Ah, 14.8~V). Fitted: identified on the
LiPo dataset~\cite{galeotti2023lipo}; literature: from~\cite{xu2016modeling}.}
\label{tb:battery_parameter}
\small
\begin{tabular}{@{}lr@{}}
\toprule
Parameter & Value \\
\midrule
\multicolumn{2}{@{}l}{\emph{\ac{ECM} (fitted)}} \\
$R_0, R_1, R_2$ at \ac{SOC}~=~0.5 & 104, 31, 33~m$\Omega$ \\
$\tau_1$, $\tau_2$ & 25, 110~s \\
$\kappa_0$, $\kappa_1$, $\kappa_2$ & 1.25, 0.89, 4.27 \\
Median validation voltage error & 12~mV \\
\midrule
\multicolumn{2}{@{}l}{\emph{Aging (fitted)}} \\
$k_{\mathrm{cyc}}$ & $7.08 \times 10^{-4}$ \\
Fade at dataset protocol & $4.8 \times 10^{-4}$/cycle \\
\midrule
\multicolumn{2}{@{}l}{\emph{Aging (literature)}} \\
$k_\sigma$, $\sigma_{\mathrm{ref}}$ & 1.04, 0.5 \\
$k_T$, $T_{\mathrm{ref}}$ & $6.93 \times 10^{-2}$, 298.15~K \\
$k_{\delta_1}$, $k_{\delta_3}$ & $1.40 \times 10^{5}$, $-1.23 \times 10^{5}$ \\
$k_{\delta_2}$ & $-0.501$ \\
$k_t$ & $4.14 \times 10^{-10}$~s$^{-1}$ \\
\bottomrule
\end{tabular}
\end{table}

\paragraph{Mission and cost} Within a mission the complexity
varies by $\pm 10\%$ (uniform, $\delta = 0.1$) from step to step around its level.
The complexity level is not known to the controllers and may change during a mission
(\emph{switch} family); the controllers estimate it on-line from the measured throughput as the
estimated complexity level $\hat c$ (Section~\ref{sec:method:estimation}).
The cost weights are $w_T = 1$, $w_E = 3/e_{\mathrm{ref}}$ with $e_{\mathrm{ref}} = 5.17$~J/m,
$w_Q = 3$, and $w_R = 5$, and the \ac{QoS} may be violated in at most $\rho = 5\%$ of the steps.
We report the return $R = -J$ (higher is better) and, in the figures, the cost $-R$.

\paragraph{Measurement realism}
All controllers observe the battery as a deployed system would. The \ac{SOC} is coulomb-counted
against the rated capacity with a per-episode gain error of up to $\pm 1\%$, the \ac{SOH} is not
observed, and the voltage and current measurements have a noise of 15~mV and 1\%, respectively.
Each episode also draws a different cell: $R_0$ is scaled by $\mathcal{U}(0.8, 1.3)$, the \ac{OCV}
is offset by up to $\pm 20$~mV, and the aging-induced \ac{OCV} change is scaled by
$\mathcal{U}(0.7, 1.3)$. The controllers only know the nominal \ac{ECM}, so the \ac{SOH}
estimator of Section~\ref{sec:method:estimation} always operates under model mismatch.

\paragraph{Scenarios}
Table~\ref{tab:scenarios} lists the 168 test scenarios, grouped into six families, which
each target one claim. Every scenario is run with 10 seeds (common across methods), i.e.,
1680 missions per method.

\begin{table}[t]
\centering
\caption{Test scenario families (168 scenarios, 10 seeds each).}
\label{tab:scenarios}
\small
\begin{tabular}{@{}llcl@{}}
\toprule
Family & Mission & \# & Varied factors \\
\midrule
Nominal        & 1~km, full charge & 9  & \ac{SOH} $\{1, 0.9, 0.8\}$ $\times$ $c$ \\
Tight          & 5~km, limited charge & 24 & \ac{SOH} $\{1, 0.8\}$ $\times$ $c$ $\times$ $k$ \\
Switch         & 3~km, full charge & 3  & $c$ redrawn every $\{10, 30, 120\}$~s \\
Mismatch       & 2~km, full charge & 15 & 5 plants $\times$ $c$ \\
Tight-mismatch & 5~km, limited charge & 90 & \ac{SOH} $\times$ $c$ $\times$ $k$ ($\geq 1$) $\times$ 5 plants \\
Degradation    & 5~km, fixed charge & 27 & \ac{SOH} $1 \to 0.8$ (step 0.025) $\times$ $c$ \\
\bottomrule
\end{tabular}
\end{table}

In the \emph{tight} families the initial \ac{SOC} is set so that the usable energy above the
reserve is $k \in \{0.9, 1.0, 1.1, 1.25\}$ times the energy of the energy-optimal mission
($k = 1.25$: driving fast is affordable; $k \approx 1$: only energy-optimal driving fits;
$k = 0.9$: the reserve must be used). The \emph{mismatch} plants differ from the nominal model
by a 50\% heavier payload, a $2\times$ or $0.7\times$ rolling/terrain term ($v^3$), or a
$1.3\times$ or $0.8\times$ computing power. In the \emph{degradation} family the mission and the
initial charge are fixed, and only the \ac{SOH} changes; the initial \ac{SOC} is chosen so that
energy-optimal driving ends exactly at the reserve at \ac{SOH}~=~80\%.

\paragraph{Compared methods}
All rule-based baselines estimate the complexity level $\hat c$ by an exponential moving average
(weight 0.3) of the complexity inverted from the observed throughput \eqref{eq:fq}, and run the
\ac{CPU} at the lowest \ac{DVFS} level that meets $q_{\mathrm{ref}}$ at the $(1-\rho)$ quantile of the
per-step complexity, $\hat c\,[1 + \delta(1 - 2\rho)]$.
\begin{itemize}
  \item \textbf{AS-MF}: constant average speed (2.5~m/s) at the maximum \ac{CPU} frequency;
        neither \ac{QoS}- nor battery-aware.
  \item \textbf{QoS-EPD}: at every step, the speed that minimizes the energy per meter
        subject to the \ac{QoS} constraint; battery-agnostic.
  \item \textbf{BA-EPD}: QoS-EPD with a charge-budget check every 10~s. It drives fast
        (3.9~m/s, or the fastest \ac{QoS}-feasible speed) if the energy for the remaining distance
        fits above the reserve with a 5\% safety factor, otherwise energy-optimally (using the
        reserve if needed), and relaxes the \ac{QoS} only if even that does not fit. It plans with
        the observed \ac{SOC} and the nominal power model, and assumes \ac{SOH}~=~1.
  \item \textbf{BA-EPD+SOH}: BA-EPD with the same on-line \ac{SOH} estimator as ours, i.e., with the
        same battery information as the proposed methods.
  \item \textbf{DARO} (proposed, Section~\ref{sec:method}) with $w_D = 16000$ unless stated
        otherwise; Section~\ref{sec:exp:longterm} varies $w_D \in \{0, 1000, 4000, 16000\}$.
\end{itemize}
DARO without the degradation cost ($w_D = 0$), i.e., Lag-RLS, and variants of it,
including one that reads the true \ac{SOC} and \ac{SOH}, are evaluated in the ablation study
(Section~\ref{sec:exp:ablation}).

\paragraph{Metrics}
Besides the return, we report the completion rate, the mission time, the load energy,
the share of steps that violate \ac{QoS}, the share of missions that end below the 5\% reserve
(reserve violation), the capacity loss per mission $\Delta\mathrm{SOH}$ (Xu model), the
distance to end of life, and the decision time. Means are given with 95\% bootstrap
confidence intervals; differences between methods are paired over scenarios and seeds.

% ---------------------------------------------------------------------
\subsection{Overall Performance}
\label{sec:exp:overall}

\begin{figure*}[t]
  \centering
  \includegraphics[width=\textwidth]{figures/exp_overall.pdf}
  \caption{Cost $-R$ per scenario family (mean over scenarios and seeds, 95\% confidence
  interval; lower is better). Bars above the plotted range are labeled with their value.}
  \label{fig:exp_overall}
\end{figure*}

\begin{table*}[t]
\centering
\caption{Return per scenario family (mean over scenarios and seeds; higher is better).
The best method in each column is in bold.}
\label{tab:return}
\small
\begin{tabular}{@{}lccccccc@{}}
\toprule
Method & Nominal & Tight & Switch & Mismatch & Tight-mism. & Degradation & All \\
\midrule
AS-MF                    & $-1.077$ & $-3.744$ & $-1.071$ & $-1.318$ & $-1.939$ & $-1.553$ & $-2.018$ \\
QoS-EPD                  & $-0.304$ & $-0.573$ & $-0.491$ & $-0.612$ & $-0.676$ & $-0.339$ & $-0.578$ \\
BA-EPD                   & $\mathbf{-0.207}$ & $-0.636$ & $-0.388$ & $-0.604$ & $-1.383$ & $-0.312$ & $-0.954$ \\
BA-EPD+SOH               & $\mathbf{-0.207}$ & $-0.538$ & $-0.388$ & $-0.604$ & $-0.811$ & $-0.280$ & $-0.628$ \\
DARO ($w_D$=16000) & $-0.250$ & $\mathbf{-0.518}$ & $\mathbf{-0.326}$ & $\mathbf{-0.536}$ & $\mathbf{-0.578}$ & $\mathbf{-0.263}$ & $\mathbf{-0.493}$ \\
\bottomrule
\end{tabular}
\end{table*}

\begin{table*}[t]
\centering
\caption{Short- and long-term metrics over all 1680 missions per method.}
\label{tab:metrics}
\small
\begin{tabular}{@{}lccccccc@{}}
\toprule
       & Completion & Time & Energy & \ac{QoS} viol. & Reserve viol. & $\Delta$\ac{SOH} & Decision \\
Method & [\%] & [min] & [kJ] & [\%] & [\%] & [ppm] & [ms] \\
\midrule
AS-MF                    & 98.2 & 29.90 & 30.39 & 0.00 & 43.8 & 32.70 & 0.00 \\
QoS-EPD                  & 100.0 & 23.27 & 26.92 & 0.97 & 5.4 & 30.05 & 0.38 \\
BA-EPD                   & 97.9 & 19.59 & 28.24 & 0.43 & 20.2 & 31.11 & 0.66 \\
BA-EPD+SOH               & 100.0 & 20.17 & 28.05 & 0.55 & 15.1 & 30.97 & 1.49 \\
DARO ($w_D$=16000) & 100.0 & 21.47 & 27.14 & 0.06 & 3.8 & 30.23 & 20.05 \\
\bottomrule
\end{tabular}
\end{table*}

\begin{figure*}[t]
  \centering
  \includegraphics[width=\textwidth]{figures/exp_tradeoff.pdf}
  \caption{Short-term trade-offs (means over all missions). (a) Mission time vs.\ load energy;
  (b) \ac{QoS} violation (share of steps) vs.\ reserve violation (share of missions). Lower left
  is better in both panels.}
  \label{fig:exp_tradeoff}
\end{figure*}

Fig.~\ref{fig:exp_overall} and Table~\ref{tab:return} report the return per family, and
Table~\ref{tab:metrics} the underlying metrics. DARO achieves the highest return over all
scenarios ($-0.493$) and in five of the six families. The strongest baselines are QoS-EPD ($-0.578$) and BA-EPD+SOH ($-0.628$). The paired
improvement of DARO is $+0.085$ (95\% CI $[+0.080, +0.091]$) over QoS-EPD and $+0.135$
($[+0.122, +0.149]$) over BA-EPD+SOH, i.e., a 15\% and 22\% lower cost; the cost is 48\% lower than
that of BA-EPD and 76\% lower than that of AS-MF. The improvement over QoS-EPD and AS-MF is
significant in all six families, and that over BA-EPD and BA-EPD+SOH in all families except
\emph{nominal}.

Fig.~\ref{fig:exp_tradeoff} shows where the gain comes from. AS-MF is the slowest (29.9~min) and
uses the most energy, since it runs the \ac{CPU} at the maximum frequency throughout. The other
baselines occupy the two ends of the time--energy trade-off: BA-EPD and BA-EPD+SOH drive fast
(19.6 and 20.2~min) but use 4.2--4.9\% more energy than QoS-EPD and end below the reserve in
15--20\% of the missions (BA-EPD fails 2.1\%), whereas QoS-EPD saves energy but is slow
(23.3~min). DARO is 7.7\% faster than QoS-EPD at 0.8\% more energy (27.14 vs.\ 26.92~kJ),
and has the fewest reserve violations (3.8\%). It also violates the \ac{QoS} in only 0.06\% of the
steps, compared with 0.4--1.0\% for the rule-based methods, whose fixed complexity allowance does
not account for the cost of a violation.

Relative to BA-EPD+SOH, the gain is largest in the charge-limited families with model mismatch.
In \emph{tight-mismatch} ($+0.233$; $+0.805$ over BA-EPD), the baselines budget the charge with the
nominal power model, whereas DARO budgets it with the \ac{RLS}-corrected model and the measured
loss factor. In \emph{tight} ($+0.020$; $+0.118$ over BA-EPD), the rule-based switching between
fast and energy-optimal driving either overspends the charge and ends below the reserve, or it
switches to energy-optimal driving earlier than necessary; the energy price $\lambda$ instead moves
continuously between both regimes. In \emph{switch} and \emph{mismatch} ($+0.062$ and $+0.069$), the
gain comes from the joint speed--frequency selection with the expected \ac{QoS} cost and the
complexity change detection. In \emph{nominal}, BA-EPD and BA-EPD+SOH are better by $0.043$: these
missions are only 1~km long (about 4.5~min), so the 90~s of speed probing at the start of
DARO take up a third of the mission, while the charge is ample and the fast mode of BA-EPD is
nearly optimal. Without probing, the cost in this family is $0.050$ lower and below that of BA-EPD
(Section~\ref{sec:exp:ablation}).

% ---------------------------------------------------------------------
\subsection{Adaptation to Battery Aging}
\label{sec:exp:aging}

\begin{figure*}[t]
  \centering
  \includegraphics[width=\textwidth]{figures/exp_aging.pdf}
  \caption{Degradation family: same 5~km mission and initial charge, aging battery.
  (a) Cost $-R$ and (b) share of missions that end below the 5\% reserve vs.\ \ac{SOH}
  (mean over complexity levels and seeds). AS-MF (cost 1.09--2.94, reserve violated from
  \ac{SOH}~=~0.9) is outside the range.}
  \label{fig:exp_aging}
\end{figure*}

Fig.~\ref{fig:exp_aging} isolates the effect of aging: the mission and the initial
\ac{SOC} are identical, and only the \ac{SOH} decreases. Down to \ac{SOH}~=~85\% the charge
suffices for all methods except AS-MF, and the cost of all methods grows only slowly with the
higher internal resistance. Below, BA-EPD, which assumes a fresh battery, keeps its fast speed
and ends below the reserve in 67\% of the missions at \ac{SOH}~=~80\%; its cost rises from
$0.26$ to $0.57$. QoS-EPD always drives energy-optimally and never violates the reserve, but it is
the slowest method (26~min) and has the highest cost down to \ac{SOH}~=~82.5\%. BA-EPD+SOH uses the
\ac{SOH} estimate, but can only switch to the energy-optimal speed for the rest of the mission; at
\ac{SOH}~=~80\% it is slow (24.1~min) and still violates the reserve in 23\% of the missions. With
the same estimate, DARO raises the energy price continuously and, through the degradation
cost, drives slightly slower from the start; its mission time grows only from 23.0 to 23.6~min, it
ends below the reserve in only 3\% of the missions at \ac{SOH}~=~80\%, and it has the lowest cost
over the whole \ac{SOH} range ($0.28$ vs.\ $0.33$ for BA-EPD+SOH at \ac{SOH}~=~80\%).

% ---------------------------------------------------------------------
\subsection{Long-Term Battery Cost}
\label{sec:exp:longterm}

\begin{figure*}[t]
  \centering
  \includegraphics[width=\textwidth]{figures/exp_longterm.pdf}
  \caption{Long-term battery cost. (a) Capacity loss per mission vs.\ mission time (means over
  all missions); the line sweeps the degradation weight $w_D \in \{0, 1000, 4000, 16000\}$
  ($w_D = 0$: without degradation cost). AS-MF (29.9~min, 32.7~ppm) is outside the range.
  (b) Distance to end of life (\ac{SOH}~=~80\%) relative to BA-EPD, per complexity level.}
  \label{fig:exp_longterm}
\end{figure*}

Fig.~\ref{fig:exp_longterm}(a) shows that the degradation weight trades mission time for battery
life. Raising $w_D$ from 0 to 16000 reduces the capacity loss from 30.44 to 30.23~ppm per mission
at a 3.4\% longer mission time, and the return still improves ($-0.493$ vs.\ $-0.498$), because
the slower driving also avoids reserve violations (3.8\% vs.\ 5.4\%). Compared with BA-EPD,
DARO extends the distance to end of life by 1.1\% and 4.1\% at low and medium complexity
and is on par ($-0.3\%$) at high complexity (Fig.~\ref{fig:exp_longterm}(b)); BA-EPD+SOH extends
it by at most 0.8\%, and AS-MF shortens it by 4.5--8.1\% because of its higher energy use.
QoS-EPD reaches the longest life ($+1.1$ to $+5.2\%$) because it always drives at the
energy-optimal speed, but it is 8\% slower than DARO and violates the \ac{QoS} in 1\% of the
steps. The long-term gain is small in absolute terms: the depth of discharge is set by the mission
length, which the controller cannot change, and only the speed-dependent part of the energy and
the mission duration are controllable. A larger $w_D$ shifts the operating point further towards
battery life, which the operator can set according to the battery replacement cost.

% ---------------------------------------------------------------------
\subsection{Ablation Study}
\label{sec:exp:ablation}

We remove or replace one component at a time and evaluate each variant on all 1680 missions with
the same scenarios and seeds, so the change in cost is paired per mission. The degradation cost is
ablated on DARO; all other components are ablated on Lag-RLS. The variants are:
\begin{itemize}
  \item \emph{w/o degradation cost}: $w_D = 0$, i.e., Lag-RLS.
  \item \emph{w/o \ac{RLS}}: the nominal power model without on-line identification.
  \item \emph{w/o change detection}: $\hat c$ is an exponential average of the inverted complexity.
  \item \emph{w/o $\lambda$ (mode switching)}: the Lagrangian price is replaced by switching between
  the fixed driving modes of BA-EPD (fast, energy-optimal, budget), on the same models and
  estimators.
  \item \emph{tier selection}: $\lambda$ is the smallest value whose plan ends above the reserve,
  instead of the value with the lowest expected cost.
  \item \emph{w/o \ac{QoS} cost}: the expected \ac{QoS} cost \eqref{eq:qos_cost} is replaced by
  the fixed complexity allowance of the baselines.
  \item \emph{w/o \ac{SOH} estimation}, \emph{w/o speed probing}, and \emph{w/o loss factor}:
  minor components, reported in the text only.
\end{itemize}
In addition, we give Lag-RLS more information than it needs: separate motor and computing power
measurements, the measured number of events, and the true \ac{SOC} and \ac{SOH}. We report the
results over all scenarios (Fig.~\ref{fig:exp_ablation}) and under battery aging
(Fig.~\ref{fig:exp_ablation_soh}).

\begin{figure}[t]
  \centering
  \includegraphics[width=\columnwidth]{figures/exp_ablation.pdf}
  \caption{Ablation: increase in cost when one main component of Lag-RLS is removed or replaced,
  relative to Lag-RLS (paired over all 1680 missions, 95\% confidence interval; positive = worse).}
  \label{fig:exp_ablation}
\end{figure}

\paragraph{Degradation cost}
Removing the degradation cost from DARO ($w_D = 0$, denoted Lag-RLS) increases the cost by
$0.005$ ($[+0.004, +0.006]$) over all scenarios, although the return does not contain the
degradation itself: Lag-RLS drives 3.3\% faster (20.8 vs.\ 21.5~min) but ends below the reserve
more often (5.4\% vs.\ 3.8\%), in particular with an aged battery (30\% vs.\ 3\% at
\ac{SOH}~=~80\%), and its capacity loss per mission is higher (30.44 vs.\ 30.23~ppm). The effect of
$w_D$ on battery life is discussed in Section~\ref{sec:exp:longterm}. Since the remaining components
are shared by both methods, we ablate them on Lag-RLS, which isolates them from the degradation
cost.

\paragraph{Components}
Fig.~\ref{fig:exp_ablation} removes one main component at a time. The \ac{RLS} power
model and the complexity change detection contribute most ($+0.081$ and $+0.071$). Without \ac{RLS},
the charge budget relies on the nominal model and fails under mismatch ($+0.163$ in
\emph{tight-mismatch}). Without change detection, the complexity estimate $\hat c$ is an exponential
average that follows the step-to-step variation, so both the \ac{DVFS} selection and the energy
forecast of the planner become noisier; the loss is largest in \emph{switch} ($+0.157$), and the
\ac{QoS} violations rise from 0.06\% to 0.50\%. Replacing the Lagrangian price by switching between fixed driving modes
(fast, energy-optimal, and budget, as in BA-EPD, but on the same models and estimators) costs
$+0.053$. Replacing the cost-based selection of $\lambda$ by the smallest $\lambda$ that fits
above the reserve (tier selection) costs $+0.038$ and more than doubles the reserve violations
(13.2\% vs.\ 5.4\%), since it cannot trade a small reserve shortfall against mission time.
Replacing the expected \ac{QoS} cost \eqref{eq:qos_cost} by the fixed complexity allowance of the
baselines costs $+0.023$ and raises the \ac{QoS} violations from 0.06\% to 0.91\%, since a fixed
allowance is either too small at high load or wastes power at low load. The remaining components
are not shown in the figure. Removing the \ac{SOH} estimation costs $+0.023$ (see the
battery-state information below). Speed probing and the loss factor have a small effect on average
($+0.008$ and $+0.002$); probing pays off in the long \emph{tight-mismatch} missions but costs
$0.050$ in the short \emph{nominal} missions.

\paragraph{Additional sensors}
We also give Lag-RLS sensors that a rover may have but our design does not require. With separate measurements of the motor and the computing power (1\%
noise each), the \ac{RLS} fits the motor and the \ac{CPU} terms independently; the return is
$0.004$ lower ($[-0.005, -0.003]$). The total power $VI$ already identifies all three coefficients
because the speed and the \ac{DVFS} level vary independently, and the remaining error of the motor
coefficients comes from the thermal transient of the motor, which a separate measurement does not
remove. With the measured number of events per batch (2\% noise) instead of the complexity inverted
from the throughput, the return is unchanged ($-0.0002$, $[-0.0006, +0.0002]$): the inversion
\eqref{eq:fq} is already accurate, and the remaining uncertainty is the step-to-step variation of
the complexity, which neither measurement can predict. Using both sensors gives $-0.003$. Hence the
proposed methods need only the battery voltage and current and the throughput.

\begin{figure*}[t]
  \centering
  \includegraphics[width=\textwidth]{figures/exp_ablation_soh.pdf}
  \caption{Ablation under aging (degradation family). (a) Cost $-R$ and (b) reserve violation
  vs.\ \ac{SOH} for DARO, Lag-RLS (without the degradation cost), and Lag-RLS with one main
  component removed or replaced, i.e., the variants of Fig.~\ref{fig:exp_ablation}.}
  \label{fig:exp_ablation_soh}
\end{figure*}

\paragraph{Ablation under aging}
Fig.~\ref{fig:exp_ablation_soh} shows how the variants of Fig.~\ref{fig:exp_ablation} behave as the
battery ages. Down to \ac{SOH}~=~85\%, none of them ends below the reserve, and the cost of each
variant rises in parallel with that of Lag-RLS. The differences appear at \ac{SOH}~=~80\%, where the
energy-optimal mission ends exactly at the reserve. Tier selection of $\lambda$ ends below the
reserve in 67\% of the missions (cost 0.454 vs.\ 0.304 for Lag-RLS): it keeps the fastest tier that
fits the estimated budget and cannot trade a small shortfall against mission time. Without change
detection, the noisier energy forecast already causes reserve violations at \ac{SOH}~=~82.5\%
(17\%) and 47\% at \ac{SOH}~=~80\%. Without the Lagrangian price or the expected \ac{QoS} cost,
Lag-RLS violates the reserve as often as with them (30\%) but at a higher cost. The degradation cost
is the component that matters most for an aged battery: DARO drives slower from the start and
ends below the reserve in only 3\% of the missions at \ac{SOH}~=~80\%. Without \ac{RLS}, Lag-RLS has
the lowest cost in this family (0.251 vs.\ 0.267), since the plant is nominal and the on-line
correction only adds estimation noise; its benefit lies in the \emph{mismatch} families.

\paragraph{Battery-state information}
The value of the \ac{SOH} estimate is separated by Lag-RLS without it and with the true battery
state (not shown in the figures). Without the estimate, Lag-RLS
keeps its speed as the battery ages and ends below the reserve in 67\% of the missions at
\ac{SOH}~=~80\%, as often as BA-EPD; over all scenarios, its return drops by $0.023$, mostly in the
charge-limited families with an aged battery. With the estimate, the reserve violations at
\ac{SOH}~=~80\% fall to 30\% (3\% for DARO). Lag-RLS with the true \ac{SOC} and \ac{SOH}
avoids them completely, but its return over all scenarios is only $0.002$ higher (reserve
violations 3.6\% vs.\ 5.4\%), which quantifies the cost of estimating the battery state from noisy
measurements of a mismatched cell.

% ---------------------------------------------------------------------
\subsection{Computational Cost}
\label{sec:exp:cost}

The mean decision time of DARO is 20.1~ms per step in our Python implementation
(19.4~ms without the degradation cost). Most of it is the joint speed--frequency search with the expected
\ac{QoS} cost (1.9~ms with the fixed allowance); the \ac{SOH} estimator accounts for about 1.4~ms
(18.1~ms with the true battery state). The decision time is two orders of magnitude below the
1~s control period, and the expected \ac{QoS} cost could be tabulated per complexity level
for an embedded implementation.
